% TOP_PROBLEM: {"id": 3063, "title": "Rota's unimodal conjecture", "category_id": 2, "category": {"id": 2, "name": "combinatorics", "display_name": "Combinatorics", "description": "Counting problems, graph theory, discrete structures.", "slug": "combinatorics", "order_index": 2, "created_at": "2026-07-31T15:26:25.671Z"}, "status": "open", "problem_number": "OPG-361", "set_id": 12}
% FIRST_POSTED: 2026-10-01
% ENTRY_KIND: statement_only
% UnsolvedMath ID 3063; source catalogue code OPG-361
% !TeX program = lualatex
% Definition and sources only; this file is not an LLM solution attempt.
\documentclass[11pt,a4paper]{article}
\usepackage[margin=25mm]{geometry}
\usepackage{fontspec}
\setmainfont{lmroman10-regular.otf}[BoldFont=lmroman10-bold.otf,
 ItalicFont=lmroman10-italic.otf,BoldItalicFont=lmroman10-bolditalic.otf]
\usepackage{amsmath,amssymb,mathtools}
\usepackage{unicode-math}
\setmathfont{latinmodern-math.otf}
\AtBeginDocument{\let\setminus\smallsetminus}
\usepackage{xurl,hyperref}
\hypersetup{colorlinks=true,urlcolor=blue}
\setcounter{secnumdepth}{0}
\setlength{\parindent}{0pt}
\setlength{\parskip}{5pt}
\setlength{\emergencystretch}{3em}
\begin{document}
\section*{Rota's unimodal conjecture}\label{3063}
{\small UnsolvedMath ID 3063 · OPG-361 · Combinatorics · OpenGarden}
\subsection{Definitions and mathematical statement}
A finite matroid may be specified by a finite set $E$ and a rank function
$r:2^E\to\mathbb Z_{\ge0}$ satisfying $0\le r(X)\le|X|$,
monotonicity $X\subseteq Y\Rightarrow r(X)\le r(Y)$, and submodularity
$r(X\cup Y)+r(X\cap Y)\le r(X)+r(Y)$. Put $R=r(E)$.
A flat, or closed set, is a subset $F\subseteq E$ such that
$r(F\cup\{e\})>r(F)$ for every $e\in E\setminus F$.
For $0\le i\le R$, let $w_i$ be the number of flats of rank $i$.

The first conjecture asserts that this sequence is unimodal: there is
an index $j\in\{0,\ldots,R\}$ such that
\[
 w_0\le\cdots\le w_j\ge w_{j+1}\ge\cdots\ge w_R.
\]
The accompanying stronger conjecture is log-concavity:
\[
 w_i^2\ge w_{i-1}w_{i+1}\qquad(1\le i<R).
\]
Both assertions quantify over every finite matroid, without any assumption
of representability over a field. For rank zero or one, the displayed
log-concavity conditions are empty.

The numbers count flats themselves, each once, even if the matroid has
loops or parallel elements. They are the Whitney numbers of the second
kind of the lattice of flats. They are not the numbers of independent
sets of each cardinality, and are not the absolute coefficients of the
characteristic polynomial. Results concerning those different sequences
do not automatically settle either assertion here.
\subsection{Short English statement}
Are the numbers of flats of successive ranks in every finite matroid unimodal, and more strongly log-concave?
\subsection{Sources}
\begin{itemize}
\item[S1] ULAM.AI and the UnsolvedMath Contributors, supplied problems.json, record 3063 (OPG-361), OpenGarden collection. Catalogue identity and source transcription; CC BY 4.0.
\url{https://huggingface.co/datasets/ulamai/UnsolvedMath}.
\item[S2] Open Problem Garden, Rota's unimodal conjecture. Original problem and discussion; the mathematical formulation below is expanded from the supplied source record.
\url{http://www.openproblemgarden.org/op/rotas_unimodal_conjecture}.
\end{itemize}
\end{document}
